\exercisesection{\S~11}

\begin{enumerate}
\def\labelenumi{\arabic{enumi})}
\tightlist
\item
  Let \(\mathfrak{g} = \mathfrak{sl}(2, k)\) . Put \(G = \operatorname{Aut}_e(\mathfrak{g})\) ; this group can be identified with \(\mathbf{PSL}_2(k)\) , cf.~Chap. VII, §3, no. 1, Remark 2.
\end{enumerate}

\begin{enumerate}
\def\labelenumi{\alph{enumi})}
\item
  Every nilpotent element of \(\mathfrak{g}\) is G-conjugate to \(\left( \begin{array}{cc}0 & \lambda \\ 0 & 0 \end{array} \right)\) for some \(\lambda \in k\) . Such an element is principal if and only if it is non-zero.
\item
  The elements \(\begin{pmatrix}0&\lambda\\0&0\end{pmatrix},\begin{pmatrix}0&\mu\\0&0\end{pmatrix}\) , where \(\lambda,\mu\in k^{*}\) , are G-conjugate if and only if \(\lambda^{-1}\mu\) is a square in k.
\item
  Every simple element of \(\mathfrak{g}\) is G-conjugate to \(\left( \begin{array}{cc}1 & 0\\ 0 & -1 \end{array} \right)\) .
\end{enumerate}

\begin{enumerate}
\def\labelenumi{\arabic{enumi})}
\setcounter{enumi}{1}
\tightlist
\item
  Let \(A = \begin{pmatrix} 0 & 1 \\ 0 & 0 \end{pmatrix}\) , \(B = \begin{pmatrix} 0 & 0 \\ 1 & 0 \end{pmatrix}\) . Then \(A\) is nilpotent, \(AB\) is not, and
\end{enumerate}

\[
[ A, [ A, B ] ] = \left( \begin{array}{c c} 0 & - 2 \\ 0 & 0 \end{array} \right).
\]

Deduce that Lemma 5 does not extend to fields of characteristic 2.

\begin{enumerate}
\def\labelenumi{\arabic{enumi})}
\setcounter{enumi}{2}
\tightlist
\item
  Let \(\mathfrak{r}\) be the radical of \(\mathfrak{g}\) , and \(\mathfrak{s} = \mathfrak{g} / \mathfrak{r}\) . Show the equivalence of:
\end{enumerate}

\begin{enumerate}
\def\labelenumi{(\roman{enumi})}
\item
  g contains no \(sl_{2}\) -triplet;
\item
  \(\mathfrak{s}\) contains no \(\mathfrak{sl}_2\) -triplet;
\item
  \(\mathfrak{s}\) is anisotropic (§10, Exerc. 6);
\item
  \(\mathfrak{s}\) contains no non-zero diagonalizable element (§3, Exerc. 10).
\end{enumerate}

(Use Prop. 2 and Exerc. 10 a) of §3.)

\begin{enumerate}
\def\labelenumi{\arabic{enumi})}
\setcounter{enumi}{3}
\tightlist
\item
  Let V be a vector space of dimension \(n \geq 2\) , \(\mathfrak{g} = \mathfrak{sl}(V)\) and \(G = \mathbf{PGL}(V)\) , identified with a group of automorphisms of g. An \(sl_{2}\) -triplet in g gives V a faithful \(\mathfrak{sl}(2, k)\) -module structure, and conversely any such structure arises from an \(sl_{2}\) -triplet; an \(sl_{2}\) -triplet is principal if and only if the corresponding \(\mathfrak{sl}(2, k)\) -module is simple; two \(sl_{2}\) -triplets are G-conjugate if and only if the corresponding \(\mathfrak{sl}(2, k)\) -modules are isomorphic. Deduce that the G-conjugacy classes of \(sl_{2}\) -triplets correspond bijectively with the families \((m_{1}, m_{2}, \ldots)\) of integers \(\geq 0\) such that
\end{enumerate}

\[
m _ {1} + 2 m _ {2} + 3 m _ {3} + \dots = n \quad \text { and } \quad m _ {1} <   n.
\]

¶ 5) Assume that g is semi-simple. Let a be a subalgebra of g, reductive in g, of the same rank as g, and containing a principal \(sl_{2}\) -triplet of g. Show that \(a = g\) .

¶6) Assume that g is absolutely simple, and denote its Coxeter number by h. Let x be a nilpotent element of g. Show that \((\operatorname{ad} x)^{2h-1} = 0\) and that \((\operatorname{ad} x)^{2h-2} \neq 0\) if and only if x is principal. (Reduce to the case in which g is split, and x is contained in the subalgebra \(n_{+}\) of Prop. 10. Repeat the proof of Prop. 10.)

¶7) Assume that g is semi-simple. Let x be a nilpotent element of g. Then x is principal if and only if x is contained in a unique Borel subalgebra of g. (Reduce to the case in which k is algebraically closed. Use Prop. 10 and Prop. 10 of §3, no. 3.)

\begin{enumerate}
\def\labelenumi{\arabic{enumi})}
\setcounter{enumi}{7}
\item
  A semi-simple Lie algebra has an \(\mathfrak{sl}_2\) -triplet if and only if it is \(\neq 0\) and has a Borel subalgebra.
\item
  Assume that g is semi-simple. Let N (resp. P) be the set of nilpotent (resp. principal nilpotent) elements of g.
\end{enumerate}

\begin{enumerate}
\def\labelenumi{\alph{enumi})}
\item
  Show that \(\mathrm{P}\) is an open subset of \(\mathbf{N}\) in the Zariski topology (use Exerc. 6).
\item
  Assume that g is splittable. Show that P is dense in N. (Use Prop. 10 and Cor. 2 of Th. 1 of §10.)
\end{enumerate}

\begin{enumerate}
\def\labelenumi{\arabic{enumi})}
\setcounter{enumi}{9}
\tightlist
\item
  Assume that \(\mathfrak{g}\) is splittable semi-simple. Let \((x,h,y)\) be an \(\mathfrak{sl}_2\) -triplet in \(\mathfrak{g}\) .
\end{enumerate}

\begin{enumerate}
\def\labelenumi{\alph{enumi})}
\item
  Show that there exists a splitting Cartan subalgebra \(\mathfrak{h}\) of \(\mathfrak{g}\) containing \(h\) . (Use Exerc. 10 b) of §3.)
\item
  Choose \(\mathfrak{h}\) as in \(a\) ). Show that \(h \in \mathfrak{h}_{\mathbf{Q}}\) and that there exists a basis B of \(\mathrm{R}(\mathfrak{g}, \mathfrak{h})\) such that \(\alpha(h) \in \{0, 1, 2\}\) for all \(\alpha \in \mathrm{B}\) (cf.~Prop. 5). The element \(x\) belongs to the subalgebra of \(\mathfrak{g}\) generated by the \(\mathfrak{g}^{\alpha}\) , \(\alpha \in \mathrm{B}\) .
\item
  Deduce from a) and b), and the Jacobson-Morozov theorem, a new proof of the fact that every nilpotent element of g is contained in a Borel subalgebra (cf.~§10, Cor. 2 of Th. 1).
\end{enumerate}

¶11) Let \((x,h,y)\) be a principal \(\mathfrak{sl}_2\) -triplet in the semi-simple Lie algebra \(\mathfrak{g}\) . Give \(\mathfrak{g}\) the \(\mathfrak{sl}(2,k)\) -module structure defined by this triplet. Show that the module thus defined is isomorphic to \(\bigoplus_{i=1}^{l} V(2k_i - 2)\) , where the \(k_i\) are the characteristic degrees of the algebra of invariant polynomial functions on \(\mathfrak{g}\) . (Reduce to the case in which \(\mathfrak{g}\) is splittable simple. Use Cor. 1 of Th. 1 of §8, no. 3, and Chap. VI, §4, Exerc. 6 c.) \(^{20}\)

¶ 12) Assume that g is semi-simple. Let \(x \in g\) and let s (resp. n) be the semi-simple (resp. nilpotent) component of x. Let \(a_{x}\) (resp. \(a_{s}\) ) be the commutant of x (resp. s) in g.

\begin{enumerate}
\def\labelenumi{\alph{enumi})}
\item
  Show that \(n\) is a nilpotent element of the semi-simple algebra \(\mathcal{D}(\mathfrak{a}_s)\) , and that the commutant of \(n\) in \(\mathfrak{a}_s\) is equal to \(\mathfrak{a}_x\) . Deduce that \(\dim \mathfrak{a}_x < \dim \mathfrak{a}_s\) if \(n \neq 0\) , i.e.~if \(x\) is not semi-simple.
\item
  Show that \(\dim \mathfrak{a}_x = \mathrm{rk}(\mathfrak{g})\) if and only if \(n\) is a principal nilpotent element of \(\mathcal{D}(\mathfrak{a}_s)\) .
\item
  Put \(\mathrm{G} = \mathrm{Aut}_{e}(\mathfrak{g})\) . Show that, for all \(\lambda \in k\) , there exists \(\sigma_{\lambda} \in G\) such that \(\sigma_{\lambda}x = s + \lambda^{2}n\) . (Show that, if \(n \neq 0\) , there exists an \(sl_{2}\) -triplet in \(a_{s}\) of which the first component is n, and deduce a homomorphism \(\varphi : \mathbf{SL}(2, k) \to \mathbf{G}\) ; take \(\sigma_{\lambda}\) to be the image under \(\varphi\) of a suitable diagonal element of \(\mathbf{SL}(2, k)\) .) Deduce that s belongs to the closure of G.x in the Zariski topology.
\end{enumerate}

\(d\) ) Show that, if \(x\) is not semi-simple, \(x\) does not belong to the closure of G.s in the Zariski topology (use the inequality \(\dim \mathfrak{a}_x < \dim \mathfrak{a}_s\) , cf.~\(a\) )).

\begin{enumerate}
\def\labelenumi{\alph{enumi})}
\setcounter{enumi}{4}
\tightlist
\item
  Assume that \(k\) is algebraically closed. Prove the equivalence of the following properties:
\end{enumerate}

\begin{enumerate}
\def\labelenumi{(\roman{enumi})}
\item
  \(x\) is semi-simple;
\item
  G.x is closed in g in the Zariski topology.
\end{enumerate}

(The implication (ii) \(\Longrightarrow\) (i) follows from c). If (i) is satisfied, and if \(x'\) is in the closure of G.x, Exerc. 15 of §8 shows that the semi-simple component \(s'\) of \(x'\) belongs to G.x, so \(x'\) belongs to the closure of G. \(s'\) ; conclude by applying d) to \(x'\) and \(s'\) .)

\(f\) ) Assume that \(k\) is algebraically closed. Let \(\mathrm{F}_x\) be the set of elements \(y\in \mathfrak{g}\) such that \(f(x) = f(y)\) for every invariant polynomial function \(f\) on \(\mathfrak{g}\) ; we have \(y\in \mathrm{F}_x\) if and only if the semi-simple component of \(y\) is G-conjugate to \(s\) (§8, Exerc. 17). Show that \(\mathrm{F}_x\) is the union of a finite number of orbits of G, and that this number is \(\leq 3^{l(x)}\) , where \(l(x)\) is the rank of \(\mathcal{D}(\mathfrak{a}_s)\) . Only one of these orbits is closed: that of \(s\) ; only one is open in \(\mathrm{F}_x\) : that consisting of the elements \(y\in \mathrm{F}_x\) such that \(\dim \mathfrak{a}_y = \operatorname {rk}(\mathfrak{g})\) .

\begin{enumerate}
\def\labelenumi{\arabic{enumi})}
\setcounter{enumi}{12}
\tightlist
\item
  Assume that \(\mathfrak{g}\) is semi-simple.
\end{enumerate}

\begin{enumerate}
\def\labelenumi{\alph{enumi})}
\item
  Let \((x, h, y)\) be a principal \(sl_{2}\) -triplet in g, and let b be the Borel subalgebra containing x (Exerc. 7). Show that b is contained in Im ad x.
\item
  Assume that k is algebraically closed. Show that, for any element z in g, there exist \(x, t \in g\) , with x principal nilpotent, such that \(z = [x, t]\) (apply a) to a Borel subalgebra containing z).
\end{enumerate}

\begin{enumerate}
\def\labelenumi{\arabic{enumi})}
\setcounter{enumi}{13}
\tightlist
\item
  Assume that \(\mathfrak{g}\) is semi-simple. Let \(\mathfrak{p}\) be a parabolic subalgebra of \(\mathfrak{g}\) , and let \(f_{1}\) and \(f_{2}\) be two homomorphisms from \(\mathfrak{g}\) to a finite dimensional Lie algebra. Show that \(f_{1}|\mathfrak{p} = f_{2}|\mathfrak{p}\) implies that \(f_{1} = f_{2}\) . (Reduce to the case in which \(\mathfrak{g}\) is split, then to the case in which \(\mathfrak{g} = \mathfrak{sl}(2,k)\) , and use Lemma 1 of no. 1.)
\end{enumerate}

¶ 15) Assume that g is semi-simple.

\begin{enumerate}
\def\labelenumi{\alph{enumi})}
\item
  Let \(x\) be a nilpotent element of \(\mathfrak{g}\) . Show that \(x\) is contained in \(\operatorname{Im}(\operatorname{ad} x)^2\) (use Prop. 2). Deduce that \((\operatorname{ad} x)^2 = 0\) implies \(x = 0\) .
\item
  Let \((x,h,y)\) be an \(\mathfrak{sl}_2\) -triplet in \(\mathfrak{g}\) , and let \(\mathfrak{s} = kx\oplus kh\oplus ky\) . Prove the equivalence of the following conditions:
\end{enumerate}

\begin{enumerate}
\def\labelenumi{(\roman{enumi})}
\item
  \(\operatorname{Im}(\operatorname{ad} x)^{2}=k.x;\)
\item
  the \(\mathfrak{s}\) -module \(\mathfrak{g} / \mathfrak{s}\) is a sum of simple modules of dimension 1 or 2;
\item
  the only eigenvalues of \(\mathrm{ad}_{\mathfrak{g}}h\) distinct from 0, 1 and -1 are 2 and -2, and their multiplicity is equal to 1.
\end{enumerate}

\begin{enumerate}
\def\labelenumi{\alph{enumi})}
\setcounter{enumi}{2}
\item
  Assume that g is splittable simple. Let h be a splitting Cartan subalgebra of g, B a basis of \(\mathrm{R}(\mathfrak{g},\mathfrak{h})\) , and \(\gamma\) the highest root of \(\mathrm{R}(\mathfrak{g},\mathfrak{h})\) relative to B. Let \((x,h,y)\) be an \(sl_{2}\) -triplet such that \(h\in h_{Q}\) and \(\alpha(h)\geq0\) for all \(\alpha\in B\) (cf.~Prop. 5). Show that conditions (i), (ii), (iii) of b) are satisfied if and only if \(h=H_{\gamma}\) , in which case \(x\in g^{\gamma}\) and \(y\in g^{-\gamma}\) .
\item
  Retain the hypotheses of c), and put \(\mathrm{G} = \mathrm{Aut}_{0}(\mathfrak{g})\) . Show that the \(sl_{2}\) -triplets satisfying (i), (ii) and (iii) are G-conjugate (use Exerc. 10). Show that, if \(x\) is a non-zero nilpotent element of \(\mathfrak{g}\) satisfying (i), the closure of \(G.x\) in the Zariski topology is equal to \(\{0\} \cup G.x\) .
\end{enumerate}

\begin{enumerate}
\def\labelenumi{\arabic{enumi})}
\setcounter{enumi}{15}
\tightlist
\item
  Let \((x,h,y)\) be an \(sl_{2}\) -triplet in g. Show that, if -2 is a square in k, the elements x-y and h are conjugate by an element of \(\operatorname{Aut}_{e}(\mathfrak{g})\) (reduce to the case in which \(\mathfrak{g}=\mathfrak{sl}(2,k)\) ). Deduce that, if g is semi-simple, x-y is semi-simple, and that it is regular if and only if h is regular.
\end{enumerate}

¶ 17) Let \((x, h, y)\) be a principal \(\mathfrak{sl}_2\) -triplet in the semi-simple algebra \(\mathfrak{g}\) . For all \(i \in \mathbf{Z}\) , denote by \(\mathfrak{g}_i\) the eigenspace of ad \(h\) relative to the eigenvalue \(i\) ; we have \(\mathfrak{g} = \bigoplus_{i \in \mathbf{Z}} \mathfrak{g}_i\) , and \(\mathfrak{g}_i = 0\) if \(i\) is odd. The direct sum \(\mathfrak{b}\) of the \(\mathfrak{g}_i\) , \(i \geq 0\) , is a Borel subalgebra of \(\mathfrak{g}\) , and \(\mathfrak{g}_0\) is a Cartan subalgebra.

\begin{enumerate}
\def\labelenumi{\alph{enumi})}
\item
  Show that, for all \(z \in \mathfrak{b}\) , the commutant of \(y + z\) in \(\mathfrak{g}\) is of dimension \(l = \operatorname{rk}(\mathfrak{g})\) .
\item
  Let I be the algebra of invariant polynomial functions on \(\mathfrak{g}\) , and \(\mathrm{P}_1, \ldots, \mathrm{P}_l\) homogeneous elements of I generating I; put \(\deg(\mathrm{P}_i) = k_i = m_i + 1\) . Show that the commutant \(\mathfrak{c}\) of \(x\) in \(\mathfrak{g}\) has a basis \(x_1, \ldots, x_l\) with \(x_i \in \mathfrak{g}_{2m_i}\) (cf.~Exerc. 11).
\item
  Let \(i \in \{1, l\}\) , and let \(J_i\) (resp. \(K_i\) ) be the set of \(j \in (1, l)\) such that \(m_j = m_i\) (resp. \(m_j < m_i\) ). Let \(f_i \in k[X_1, \ldots, X_l]\) be the polynomial such that
\end{enumerate}

\[
f _ {i} \left(a _ {1}, \dots , a _ {l}\right) = \mathrm{P} _ {i} \left(y + \sum_ {j = 1} ^ {j = l} a _ {j} x _ {j}\right) \quad \text { for } \left(a _ {j}\right) \in k ^ {l}.
\]

Show that \(f_{i}\) is the sum of a linear form \(\mathrm{L}_i\) in the \(\mathrm{X}_j\) , \(j \in \mathrm{J}_i\) , and a polynomial in the \(\mathrm{X}_j\) , \(j \in \mathrm{K}_i\) . (If \(t \in k^*\) , the automorphism of \(\mathfrak{g}\) equal to \(t^i\) on \(\mathfrak{g}_i\) belongs to \(\operatorname{Aut}_0(\mathfrak{g})\) and takes \(y\) to \(t^{-2}y\) and \(x_j\) to \(t^{2m_j}x_j\) . Use the invariance of \(\mathrm{P}_i\) under this automorphism.)

\(d)\) Let \(\mathrm{P}\) be the map from \(\mathfrak{g}\) to \(k^l\) defined by the \(\mathrm{P}_i\) . If \(z \in \mathfrak{g}\) , denote by \(\mathrm{D}_z\mathrm{P} : \mathfrak{g} \to k^l\) the tangent linear map of \(\mathrm{P}\) at \(z\) (Chap. VII, App. I, no. 2).

Show that \(\mathfrak{c}\cap\operatorname{Im}\operatorname{ad}(y-x)=0\) (decompose g into a direct sum of simple submodules relative to the subalgebra generated by the given \(sl_{2}\) -triplet). Deduce that the restriction of \(D_{y-x}P\) to c is an isomorphism from c to \(k^{l}\) (use the preceding exercise and Exerc. 18 a) of §8). Show, by using this result, that the determinant of the linear forms \(L_{i}\) defined in c) is \(\neq0\) , and deduce the following results:

\(d_{1})\) the polynomials \(f_{1},\ldots ,f_{l}\) are algebraically independent and generate \(k[\mathrm{X}_1,\dots ,\mathrm{X}_l]\) ;

\(d_{2})\) the map \(z \mapsto \mathrm{P}(y + z)\) from \(\mathfrak{c}\) to \(k^{l}\) is bijective polynomial, and the inverse map is polynomial;

\(d_{3})\) for all \(z\in y + \mathfrak{c}\) , the linear map \(\mathrm{D}_z\mathrm{P}|\mathfrak{c}\) is of rank \(l\) .

In particular, the map \(\mathrm{P}:\mathfrak{g}\to k^l\) is surjective.

\begin{enumerate}
\def\labelenumi{\alph{enumi})}
\setcounter{enumi}{4}
\tightlist
\item
  If \(k\) is algebraically closed, every element of \(\mathfrak{g}\) whose commutant is of dimension \(l\) is conjugate under \(\mathrm{Aut}_0(\mathfrak{g})\) to a unique element of \(y + \mathfrak{c}\) .
\end{enumerate}

\(f\) ) Give an example of a simple Lie algebra with no principal \(\mathfrak{sl}_2\) -triplet, for which the map \(\mathrm{P}\) is not surjective (take \(k = \mathbf{R}\) and \(l = 1\) ).

